\documentclass[11pt]{amsart}

\usepackage[margin=1.15in]{geometry}
\usepackage{amsmath,amssymb,amsthm,mathtools}
\usepackage{enumitem}
\usepackage{microtype}
\usepackage{hyperref}
\usepackage[nameinlink,capitalize]{cleveref}

\hypersetup{
  colorlinks=true,
  linkcolor=blue,
  citecolor=blue,
  urlcolor=blue
}

\theoremstyle{plain}
\newtheorem{theorem}{Theorem}[section]
\newtheorem{proposition}[theorem]{Proposition}
\newtheorem{lemma}[theorem]{Lemma}
\newtheorem{corollary}[theorem]{Corollary}

\theoremstyle{definition}
\newtheorem{definition}[theorem]{Definition}
\newtheorem{assumption}[theorem]{Assumption}

\theoremstyle{remark}
\newtheorem{remark}[theorem]{Remark}

\numberwithin{equation}{section}

\newcommand{\R}{\mathbb R}
\newcommand{\N}{\mathbb N}
\newcommand{\eps}{\varepsilon}
\newcommand{\loc}{\mathrm{loc}}
\newcommand{\supp}{\operatorname{supp}}
\newcommand{\divergence}{\operatorname{div}}
\newcommand{\dist}{\operatorname{dist}}
\newcommand{\esssup}{\operatorname*{ess\,sup}}
\newcommand{\avg}[1]{\left(#1\right)}
\newcommand{\B}{B}
\newcommand{\Q}{Q}
\newcommand{\calN}{\mathcal N}
\newcommand{\calE}{\mathcal E}
\newcommand{\calI}{\mathcal I}

\title[Type I blow-up limits for Navier--Stokes]{Type I Blow-up Limits for the Three-dimensional Navier--Stokes Equations and the Centered Ancient Equation}

\author{}
\date{}

\subjclass[2020]{35Q30, 35B44, 76D05}
\keywords{Navier--Stokes equations, Type I singularities, ancient solutions, suitable weak solutions, self-similar variables}

\begin{document}

\begin{abstract}
We give a detailed reduction from a Type I singular point of a suitable
Leray--Hopf solution of the three-dimensional incompressible Navier--Stokes
equations to a bounded ancient suitable weak solution in centered self-similar
variables.  The construction records the precise scaling, pressure gauges,
local compactness, local energy stability, and the pressure-normalized
Caffarelli--Kohn--Nirenberg concentration which prevents the limit from
vanishing.  The compactness and nonvanishing steps use the
Caffarelli--Kohn--Nirenberg--Seregin blow-up compactness mechanism in the
precise form recorded below.  The resulting renormalized ancient field
solves
\[
   \partial_\tau V+(V\cdot\nabla)V+\nabla \Pi
   =\Delta V-\frac12(V+y\cdot\nabla V),\qquad \divergence V=0,
\]
on \(\R^3\times\R\), and is uniformly bounded in \(L^\infty_{y,\tau}\).  We
also observe that this centered equation has no non-trivial positive scaling
symmetry and is not invariant under constant spatial translations.  Parabolic
rescalings of the unrenormalized ancient solution appear only as translations
in the logarithmic time variable.  Thus centered Type I ancient limits are not
organized by an additional dilation orbit in the centered variables.
\end{abstract}

\maketitle

\section{Introduction}

Let \(u:\R^3\times[0,T)\to\R^3\) solve the incompressible Navier--Stokes
equations
\begin{equation}\label{eq:NS}
   \partial_t u+(u\cdot\nabla)u+\nabla p=\Delta u,\qquad
   \divergence u=0.
\end{equation}
The local regularity theory of Caffarelli--Kohn--Nirenberg
\cite{CKN1982} shows that singular points are detected by scale-invariant
concentration of velocity and pressure.  A finite-time singular point
\(z_*=(x_*,T)\) is said to be of Type I if the scale-invariant \(L^\infty\)
quantity remains bounded,
\begin{equation}\label{eq:type-I-bound-intro}
   \limsup_{t\uparrow T}\sqrt{T-t}\,
   \|u(t)\|_{L^\infty(\R^3)}<\infty.
\end{equation}
For the endpoint compactness argument one also needs the scale-invariant
energy and pressure quantities at cylinders ending at \(z_*\) to remain
finite.  We record this as the hypothesis
\(\calI_*(u,p;z_*)<\infty\).  This is the local compactness formulation of
Type I used in the partial regularity literature; it follows from the standard
critical Morrey or weak Serrin Type I hypotheses, and it is the condition used
below to pass to a suitable limit up to the terminal slice.
For the relation between local Type I hypotheses, scale-invariant energy
control, and bounded ancient limits, see Albritton--Barker
\cite[Theorem 1.1]{AlbrittonBarker2019}.
The purpose of this paper is to isolate the canonical ancient object associated
with such a singular point.

We use the pressure-normalized CKN quantities, their parabolic scaling, the
epsilon-regularity/no-concentration implication, and the local Type I/Type II
rate dichotomy in the form stated and proved below.  Thus the compactness
argument is self-contained at the level needed here: the only external inputs
are the standard Caffarelli--Kohn--Nirenberg regularity criterion and the
classical Seregin and Albritton--Barker compactness theorems cited explicitly.
The additional work in the present paper is the Type I compactness up to
ancient time, the pressure gauges needed for the whole-space limit, the
nonvanishing normalization of the selected limit, and the centered
self-similar equation.

The parabolic scaling around \(z_*\) is
\begin{equation}\label{eq:basic-rescaling-intro}
   u^{(\lambda)}(x,t)
   =
   \lambda u(x_*+\lambda x,T+\lambda^2t),\qquad
   p^{(\lambda)}(x,t)
   =
   \lambda^2 p(x_*+\lambda x,T+\lambda^2t).
\end{equation}
As \(\lambda\downarrow0\), the time intervals
\(( -T/\lambda^2,0)\) exhaust \((-\infty,0)\).  The Type I bound becomes
\[
   \|u^{(\lambda)}(t)\|_{L^\infty(\R^3)}
   \lesssim (-t)^{-1/2},
   \qquad t<0,
\]
uniformly on compact subintervals of \((-\infty,0)\).  This gives local
compactness away from \(t=0\).  After selecting scales carefully, one obtains a
non-zero ancient suitable weak limit \(U\).

The centered self-similar variables associated with the terminal time are
\begin{equation}\label{eq:self-sim-vars-intro}
   y=\frac{x}{\sqrt{-t}},\qquad
   \tau=-\log(-t),\qquad
   V(y,\tau)=\sqrt{-t}\,U(y\sqrt{-t},t).
\end{equation}
Then \(V\) solves the autonomous equation
\begin{equation}\label{eq:centered-intro}
   \partial_\tau V+(V\cdot\nabla)V+\nabla\Pi
   =
   \Delta V-\frac12(V+y\cdot\nabla V),
   \qquad \divergence V=0.
\end{equation}
The drift term in \eqref{eq:centered-intro} is the signature of the Type I
scaling.  It fixes a scale and center.  Consequently, unlike the original
Navier--Stokes equations, \eqref{eq:centered-intro} has no non-trivial
amplitude-space scaling symmetry and no constant translation symmetry in the
\(y\)-variable.  The remaining symmetry is translation in \(\tau\), which is
the image of parabolic rescaling before the centered change of variables.
Consequently, after the centered change of variables has been made, Type I
ancient dynamics are not generated by further same-center parabolic dilations.

\subsection{Main result}

The following theorem is the main result of the paper.  The non-triviality
assertion is stated in the pressure-normalized CKN form supplied by the local
concentration theorem at singular points.

\begin{theorem}[Type I blow-up limit and centered ancient equation]
\label{thm:main}
Let \((u,p)\) be a suitable Leray--Hopf solution of \eqref{eq:NS} on
\(\R^3\times(0,T)\).  Assume \(z_*=(x_*,T)\) is a singular point and that
\eqref{eq:type-I-bound-intro} holds.  Assume in addition that the centered
endpoint Type I compactness quantity is finite:
\begin{equation}\label{eq:centered-endpoint-I-main}
   \calI_*(u,p;z_*)<\infty .
\end{equation}
Here \(\calI_*\) is defined in \eqref{eq:centered-I-def}.
Then there are scales
\(\lambda_k\downarrow0\), functions of time \(a_k(t)\), and an ancient
suitable weak solution \((U,P)\) on \(\R^3\times(-\infty,0)\) such that, with
\[
   u_k(x,t)=\lambda_k u(x_*+\lambda_k x,T+\lambda_k^2t),
   \qquad
   \widetilde p_k(x,t)=\lambda_k^2p(x_*+\lambda_k x,T+\lambda_k^2t),
   \qquad
   \pi_k=\widetilde p_k-a_k(t),
\]
one has, after passing to a subsequence,
\[
   u_k\to U \quad\hbox{strongly in } L^q_{\loc}(\R^3\times(-\infty,0))
   \quad\hbox{for every }1\le q<\infty,
\]
and
\[
   \pi_k\rightharpoonup P
   \quad\hbox{weakly in }L^{3/2}_{\loc}(\R^3\times(-\infty,0)).
\]
The limit satisfies
\begin{equation}\label{eq:ancient-type-I-bound}
   \|U(t)\|_{L^\infty(\R^3)}
   \le \frac{M}{\sqrt{-t}},
   \qquad t<0,
\end{equation}
for a finite constant \(M\) depending only on the Type I bound.  Moreover the
scales may be chosen so that \(U\not\equiv0\).  The renormalized field \(V\)
defined by \eqref{eq:self-sim-vars-intro} satisfies
\[
   \|V\|_{L^\infty(\R^3\times\R)}\le M
\]
and is a smooth solution of \eqref{eq:centered-intro} on \(\R^3\times\R\).
The non-triviality is quantified as follows: there are a compact cylinder
\[
   K\Subset \R^3\times\R,
\]
a function \(b(\tau)\), and \(\eta_0>0\) such that
\begin{equation}\label{eq:main-nonzero-renorm}
   \iint_K\left(|V|^3+|\Pi-b(\tau)|^{3/2}\right)\,dy\,d\tau\ge\eta_0 .
\end{equation}
\end{theorem}

\begin{proposition}[No additional dilation or translation symmetry]
\label{prop:no-cascade}
The centered equation \eqref{eq:centered-intro} is autonomous in \(\tau\), but
it is not invariant under the positive amplitude-space scaling
\[
   V(y,\tau)\mapsto \alpha V(\alpha y,\alpha^2\tau)
\]
unless \(\alpha=1\).  It is also not invariant under translations
\[
   V(y,\tau)\mapsto V(y-a,\tau)
\]
unless \(a=0\).  Hence the centered Type I equation has no intrinsic
same-center parabolic dilation symmetry beyond translation of the logarithmic
time variable.
\end{proposition}

\subsection{Organization}

Section \ref{sec:prelim} fixes notation and recalls suitable weak solutions,
scale-invariant quantities, and epsilon regularity.  Section
\ref{sec:bounds} proves uniform bounds for the rescaled sequence.  Section
\ref{sec:compactness} passes to an ancient suitable weak limit.  Section
\ref{sec:nontrivial} proves the non-triviality normalization.  Section
\ref{sec:renorm} derives the centered equation.  Section \ref{sec:symmetry}
proves \cref{prop:no-cascade}.

\subsection{Proof bookkeeping}

The proof has three logically separate layers.  First, the Type I bound gives
local \(L^\infty\) compactness on every cylinder separated from the terminal
time.  Second, the finite endpoint compactness quantity
\(\calI_*(u,p;z_*)\) gives compactness up to rescaled cylinders ending at
\((0,0)\), after subtracting time-dependent pressure gauges.  Third, the
positive pressure-normalized CKN concentration at the singular point chooses
the scales so that the limiting ancient solution is not zero.  None of these
three layers implies the others.  In particular, the Type I bound alone gives
boundedness of possible limits but not non-triviality, while local
concentration alone gives non-triviality but not the global-in-time
\(L^\infty\) ancient bound.

\section{Preliminaries}\label{sec:prelim}

We write
\[
   Q_r(z_0)=B_r(x_0)\times(t_0-r^2,t_0),
   \qquad z_0=(x_0,t_0).
\]
When \(z_0=(0,0)\) we write \(Q_r=Q_r(0,0)\).  The spatial average of a
function \(f\) over \(B_r(x_0)\) is denoted by
\[
   (f)_{B_r(x_0)}(t)=\frac1{|B_r|}\int_{B_r(x_0)}f(x,t)\,dx.
\]

\begin{definition}[Terminal singular point]\label{def:terminal-singular}
Let \((u,p)\) be a suitable weak solution on \(\R^3\times(0,T)\).  A terminal
point \(z_*=(x_*,T)\) is called regular if there are \(r>0\) and
\(K<\infty\) such that
\[
   \esssup_{Q_r(z_*)}|u|\le K.
\]
It is called singular otherwise.
\end{definition}

\begin{definition}[Suitable Leray--Hopf solution]\label{def:suitable}
Let \(\Omega\subset \R^3\times\R\) be open.  A pair \((u,p)\) is a suitable
weak solution of \eqref{eq:NS} in \(\Omega\) if
\[
   u\in L^\infty_tL^2_x(\Omega')\cap L^2_t\dot H^1_x(\Omega'),
   \qquad
   p\in L^{3/2}(\Omega')
\]
for every compact cylinder \(\Omega'\Subset\Omega\), if \((u,p)\) solves
\eqref{eq:NS} distributionally, \(\divergence u=0\), and if the local energy
inequality
\begin{align}\label{eq:lei}
   &\int_{\R^3}|u(x,t_2)|^2\phi(x,t_2)\,dx
   +2\int_{t_1}^{t_2}\int_{\R^3}|\nabla u|^2\phi\,dx\,dt  \notag\\
   &\le
   \int_{\R^3}|u(x,t_1)|^2\phi(x,t_1)\,dx
   +\int_{t_1}^{t_2}\int_{\R^3}|u|^2(\partial_t\phi+\Delta\phi)\,dx\,dt
   \notag\\
   &\quad
   +\int_{t_1}^{t_2}\int_{\R^3}(|u|^2+2p)u\cdot\nabla\phi\,dx\,dt
\end{align}
holds for every non-negative \(\phi\in C_c^\infty(\Omega)\) and almost every
\(t_1<t_2\) for which the displayed terms are defined.
\end{definition}

This is the Caffarelli--Kohn--Nirenberg notion of suitability
\cite{CKN1982}; the global finite-energy background is Leray's
Leray--Hopf class \cite{Leray1934}.

We use the standard scale-invariant quantities
\begin{equation}\label{eq:C-D-def}
   A(u;z_0,r)=r^{-1}\esssup_{t_0-r^2<t<t_0}
   \int_{B_r(x_0)}|u(x,t)|^2\,dx,
\end{equation}
\begin{equation}\label{eq:C-D-def-2}
   C(u;z_0,r)=r^{-2}\iint_{Q_r(z_0)}|u|^3\,dx\,dt,
   \qquad
   D(p;z_0,r)=r^{-2}\iint_{Q_r(z_0)}
   |p-(p)_{B_r(x_0)}(t)|^{3/2}\,dx\,dt.
\end{equation}
\begin{equation}\label{eq:E-I-def}
   E(u;z_0,r)=r^{-1}\iint_{Q_r(z_0)}|\nabla u|^2\,dx\,dt,
\end{equation}
\begin{equation}\label{eq:I-def}
   \calI(u,p;z_0,r)=
   A(u;z_0,r)+C(u;z_0,r)+D(p;z_0,r)+E(u;z_0,r).
\end{equation}
For a terminal point \(z_*=(x_*,T)\), we write
\begin{equation}\label{eq:centered-I-def}
   \calI_*(u,p;z_*)=
   \sup_{0<r<r_*}\calI(u,p;z_*,r),
\end{equation}
where \(r_*>0\) is fixed sufficiently small that the corresponding cylinders
lie in the spacetime domain under consideration.  Finiteness of
\(\calI_*(u,p;z_*)\) is independent of decreasing \(r_*\).

\begin{theorem}[Epsilon regularity]\label{thm:eps-reg}
There exists \(\eps_*>0\) with the following property.  If \((u,p)\) is a
suitable weak solution in \(Q_r(z_0)\) and
\[
   C(u;z_0,r)+D(p;z_0,r)\le \eps_*,
\]
then \(u\) is bounded, and hence smooth, in \(Q_{r/2}(z_0)\).
\end{theorem}

This is the Caffarelli--Kohn--Nirenberg regularity criterion
\cite{CKN1982}; see also Lin's direct proof \cite{Lin1998}.  We shall use only
its contrapositive.

\begin{lemma}[Scaling]\label{lem:scaling}
Let \((u,p)\) be a suitable weak solution in a neighborhood of
\((x_*,T)\).  For \(\lambda>0\) set
\[
   u^\lambda(x,t)=\lambda u(x_*+\lambda x,T+\lambda^2t),
   \qquad
   p^\lambda(x,t)=\lambda^2 p(x_*+\lambda x,T+\lambda^2t).
\]
Then \((u^\lambda,p^\lambda)\) is a suitable weak solution on the rescaled
domain.  Moreover \(A,C,D,E\) in
\eqref{eq:C-D-def}--\eqref{eq:E-I-def} are invariant under
this scaling.
\end{lemma}

\begin{proof}
Suitability of the rescaled pair follows by pulling back the distributional
equation and testing the local energy inequality against the corresponding
nonnegative rescaled cutoffs.  The pressure gauge is preserved because
subtracting a function of time has zero spatial gradient and cancels from
pressure oscillations.  The quantities \(C\) and \(D\) are invariant by the
same parabolic Jacobian calculation used below for \(A\) and \(E\).  With
\(X=x_*+\lambda x\) and
\(S=T+\lambda^2t\),
\begin{align*}
   \int_{B_r}|u^\lambda(x,t)|^2\,dx
   &=
   \lambda^{-1}\int_{B_{\lambda r}(x_*)}|u(X,S)|^2\,dX,\\
   \iint_{Q_r}|\nabla u^\lambda|^2\,dx\,dt
   &=
   \lambda^{-1}\iint_{Q_{\lambda r}(z_*)}|\nabla u|^2\,dX\,dS .
\end{align*}
Multiplication by \(r^{-1}\) on the rescaled cylinder is therefore the same as
multiplication by \((\lambda r)^{-1}\) on the original cylinder.  Hence
\(A\) and \(E\) are invariant as well.
\end{proof}

\begin{lemma}[Pressure gauges do not change the equations]\label{lem:gauges-harmless}
Let \((v,q)\) be a suitable weak solution in a cylinder \(Q\), and let
\(\alpha(t)\in L^{3/2}_{\mathrm{loc}}\) depend only on time.  Then
\((v,q-\alpha(t))\) satisfies the same distributional Navier--Stokes equation,
the same divergence condition, and the same local energy inequality.
\end{lemma}

\begin{proof}
In the distributional equation, replacing \(q\) by \(q-\alpha(t)\) adds
\(-\nabla\alpha(t)\), which is zero because \(\alpha\) has no spatial
dependence.  The divergence equation is independent of the pressure.  In the
local energy inequality the only new term is
\[
   -2\int \alpha(t)v(x,t)\cdot\nabla\phi(x,t)\,dx\,dt .
\]
For almost every \(t\), since \(\operatorname{div}v(\cdot,t)=0\) in
distributions and \(\phi(\cdot,t)\) is compactly supported,
\[
   \int v(x,t)\cdot\nabla\phi(x,t)\,dx=0 .
\]
Thus the additional term vanishes.  Therefore all compactness arguments may
choose pressure representatives modulo functions of time without altering
suitability.
\end{proof}

\section{Uniform estimates for the rescaled sequence}\label{sec:bounds}

Let \(M<\infty\) be chosen so that, after decreasing \(T_0<T\) if necessary,
\begin{equation}\label{eq:type-I-bound}
   \|u(t)\|_{L^\infty(\R^3)}
   \le \frac{M}{\sqrt{T-t}},
   \qquad T_0<t<T.
\end{equation}
For a sequence \(\lambda_k\downarrow0\), define
\begin{equation}\label{eq:rescaled-sequence}
   u_k(x,t)=\lambda_k u(x_*+\lambda_k x,T+\lambda_k^2t),
   \qquad
   p_k(x,t)=\lambda_k^2p(x_*+\lambda_k x,T+\lambda_k^2t).
\end{equation}
Then \(u_k\) is defined on
\[
   \R^3\times(-T/\lambda_k^2,0),
\]
and these domains exhaust \(\R^3\times(-\infty,0)\).

\begin{lemma}[Inherited Type I bound]\label{lem:inherited-bound}
For every compact interval \(I\Subset(-\infty,0)\),
\[
   \sup_k\|u_k\|_{L^\infty(\R^3\times I)}<\infty.
\]
More precisely, for every fixed \(t<0\) and all sufficiently large \(k\),
\[
   \|u_k(t)\|_{L^\infty(\R^3)}
   \le \frac{M}{\sqrt{-t}}.
\]
\end{lemma}

\begin{proof}
For \(t<0\),
\[
   T-(T+\lambda_k^2t)=\lambda_k^2(-t).
\]
Thus \eqref{eq:type-I-bound} gives, for all large \(k\),
\[
   \|u_k(t)\|_\infty
   =
   \lambda_k \|u(T+\lambda_k^2t)\|_\infty
   \le
   \lambda_k\frac{M}{\sqrt{\lambda_k^2(-t)}}
   =
   \frac{M}{\sqrt{-t}}.
\]
The interval statement follows by taking the supremum over \(t\in I\).
\end{proof}

\begin{lemma}[Local energy and dissipation]\label{lem:energy-diss}
Let
\[
   Q=B_R\times[-S,-\sigma]\Subset \R^3\times(-\infty,0),
   \qquad 0<\sigma<S<\infty.
\]
Then
\[
   \sup_k\|u_k\|_{L^\infty_tL^2_x(Q)}
   +\sup_k\|\nabla u_k\|_{L^2(Q)}
   <\infty .
\]
\end{lemma}

\begin{proof}
The \(L^\infty_tL^2_x\) estimate is a direct consequence of
\cref{lem:inherited-bound}:
\[
   \|u_k(t)\|_{L^2(B_R)}
   \le |B_R|^{1/2}\frac{M}{\sqrt{\sigma}},
   \qquad t\in[-S,-\sigma].
\]
To estimate \(\nabla u_k\), choose
\(\phi\in C_c^\infty(B_{2R}\times[-2S,-\sigma/2])\), \(0\le\phi\le1\),
with \(\phi\equiv1\) on \(Q\).  Apply the local energy inequality to
\((u_k,p_k-a_k(t))\), where \(a_k(t)=(p_k)_{B_1}(t)\) is the fixed reference
gauge from \cref{lem:pressure}.  The estimate in \cref{lem:pressure} is used
on the larger cylinder \(B_{2R}\times[-2S,-\sigma/2]\).  Subtracting this
function of time from the pressure does not change the local energy inequality
by \cref{lem:gauges-harmless}.  For a.e. admissible pair of times
\(-2S<t_1<t_2<-\sigma/2\), and then by monotone approximation of the temporal
cutoff to the endpoints of the support of \(\phi\), the inequality gives
\[
   2\iint |\nabla u_k|^2\phi
   \le
   \iint |u_k|^2|\partial_t\phi+\Delta\phi|
   +\iint (|u_k|^2+2|p_k-a_k(t)|)|u_k||\nabla\phi|.
\]
The first term is bounded by the inherited \(L^\infty\) bound and the finite
measure of \(B_{2R}\times[-2S,-\sigma/2]\).  For the cubic velocity term,
\[
   \iint |u_k|^3|\nabla\phi|
   \le
   \|\nabla\phi\|_\infty
   |B_{2R}\times[-2S,-\sigma/2]|
   \|u_k\|_{L^\infty(B_{2R}\times[-2S,-\sigma/2])}^3 .
\]
For the pressure term, Hölder's inequality with exponents \(3/2\) and \(3\)
gives
\[
   \iint |p_k-a_k(t)|\,|u_k|\,|\nabla\phi|
   \le
   \|p_k-a_k(t)\|_{L^{3/2}}
   \|u_k\nabla\phi\|_{L^3},
\]
and both factors are uniformly bounded by \cref{lem:pressure} and
\cref{lem:inherited-bound}.  Absorbing no term is necessary, since the
dissipation appears with a favorable sign.  Because \(\phi\equiv1\) on \(Q\),
this yields
\[
   \iint_Q|\nabla u_k|^2\,dx\,dt\le C(R,S,\sigma,M).
\]
\end{proof}

\begin{lemma}[Pressure gauges and local pressure bounds]\label{lem:pressure}
After fixing the whole-space Leray pressure representative below, set
\[
   a_k(t)=(p_k)_{B_1}(t)
\]
whenever the rescaled solution is defined.  Then for every
\(Q=B_R\times[-S,-\sigma]\Subset\R^3\times(-\infty,0)\),
\[
   \sup_k\|p_k-a_k(t)\|_{L^{3/2}(Q)}<\infty.
\]
\end{lemma}

\begin{proof}
For all sufficiently large \(k\), the original times
\(T+\lambda_k^2t\), \(t\in[-S,-\sigma]\), lie in the interval
\((T_0,T)\), where the Type I bound \eqref{eq:type-I-bound} holds.  On those
time slices \(u(t)\in L^2(\R^3)\cap L^\infty(\R^3)\), hence
\[
   u_i(t)u_j(t)\in L^1(\R^3)\cap L^{3/2}(\R^3)
\]
by interpolation.  We use the whole-space Leray pressure representative
\[
   p^L=\mathcal R_i\mathcal R_j(u_i u_j),
\]
modulo functions of time; this is the pressure reconstruction from
\(-\Delta p=\partial_i\partial_j(u_i u_j)\) and the Calderon--Zygmund theory
of Riesz transforms \cite{Stein1970,Sohr2001}.  For whole-space Leray--Hopf
solutions the given pressure and \(p^L\) differ only by a function of time.
Indeed, both pressures give the same gradient part in the distributional
momentum equation, equivalently the same distributional Laplacian
\[
   -\Delta p=\partial_i\partial_j(u_i u_j),
\]
and the de Rham theorem applied to the difference of the two weak formulations
gives \(\nabla(p-p^L)=0\) in \(\mathcal D'(\R^3)\) for a.e. time.  Thus
\(p-p^L\) is spatially constant for a.e. time.  Replacing the original
pressure by \(p^L\) is therefore only a time-dependent gauge change, which
does not affect suitability by \cref{lem:gauges-harmless}.  The rescaled
pressure has the same representative with \(u_k\), again modulo a
time-dependent function.  In the rest of this lemma \(p_k\) denotes one such
representative; after that representative is fixed, the reference gauge
\(a_k(t)=(p_k)_{B_1}(t)\) is fixed as well.

Fix a smooth cutoff \(\chi_R\in C_c^\infty(B_{8R})\) with
\(\chi_R\equiv1\) on \(B_{4R}\).  For \(t\in[-S,-\sigma]\), decompose in
\(B_{2R}\)
\[
   p_k=q_k+h_k,
   \qquad
   q_k=\mathcal R_i\mathcal R_j\bigl(\chi_R u_{k,i}u_{k,j}\bigr).
\]
More precisely, the representative may be chosen, by adding a function of
\(t\), so that the difference \(h_k=p_k-q_k\) satisfies
\(-\Delta h_k=0\) in \(B_{4R}\) for a.e. \(t\), because the localized source
coincides with \(u_{k,i}u_{k,j}\) on \(B_{4R}\).  This choice is made before
the reference gauge \(a_k(t)=(p_k)_{B_1}(t)\) is evaluated, so the subsequent
comparison of \(a_k\) with \(c_{k,R}\) uses a single pressure representative.
Calderon--Zygmund boundedness on \(L^{3/2}\) gives
\[
   \|q_k(t)\|_{L^{3/2}(B_{2R})}
   \le C\|\chi_R u_k(t)\otimes u_k(t)\|_{L^{3/2}(\R^3)}
   \le C\|u_k(t)\|_{L^3(B_{8R})}^2
   \le C(R,\sigma,M).
\]
It remains first to control the harmonic far field modulo a constant.  Let
\[
   c_{k,R}(t)=h_k(0,t).
\]
The value at the origin is defined by the smooth harmonic representative of
\(h_k(\cdot,t)\) in \(B_{4R}\); measurability in \(t\) follows by testing the
harmonic representative against a fixed mollifier supported in a small ball.
Let \(K_{ij}\) denote the Calderon--Zygmund kernel of
\(\mathcal R_i\mathcal R_j\).  For \(x\in B_R\), the kernel estimate
\[
   |K_{ij}(x-z)-K_{ij}(-z)|\le C R |z|^{-4},
   \qquad |z|>4R,
\]
implies
\[
   |h_k(x,t)-c_{k,R}(t)|
   \le
   C R\sum_{m=2}^\infty (2^mR)^{-4}
   \int_{2^mR<|z|<2^{m+1}R}|1-\chi_R(z)|\,|u_k(z,t)|^2\,dz .
\]
By \cref{lem:inherited-bound},
\[
   \int_{2^mR<|z|<2^{m+1}R}|u_k(z,t)|^2\,dz
   \le C M^2\sigma^{-1}(2^mR)^3 .
\]
Consequently
\[
   \|h_k(t)-c_{k,R}(t)\|_{L^\infty(B_R)}
   \le
   C M^2\sigma^{-1}
   R\sum_{m=2}^\infty (2^mR)^{-1}
   \le C(M,\sigma).
\]
Combining the near-field and far-field estimates and integrating over the
finite interval \([-S,-\sigma]\) gives
\[
   \sup_k\|p_k-c_{k,R}(t)\|_{L^{3/2}(B_R\times[-S,-\sigma])}<\infty .
\]

We now pass from the local gauge \(c_{k,R}\) to the single reference gauge
\(a_k(t)=(p_k)_{B_1}(t)\).  If \(R<1\), enlarge \(R\) to \(1\) in the preceding
estimate; hence assume \(R\ge1\).  Since \(B_1\subset B_R\),
Jensen's inequality gives, for a.e. \(t\in[-S,-\sigma]\),
\[
   |a_k(t)-c_{k,R}(t)|^{3/2}
   =
   \left|(p_k-c_{k,R}(t))_{B_1}(t)\right|^{3/2}
   \le
   \frac1{|B_1|}\int_{B_1}|p_k-c_{k,R}(t)|^{3/2}\,dx .
\]
Integrating in time shows that \(a_k-c_{k,R}\) is uniformly bounded in
\(L^{3/2}([-S,-\sigma])\).  Therefore
\[
   \|p_k-a_k(t)\|_{L^{3/2}(B_R\times[-S,-\sigma])}
   \le
   \|p_k-c_{k,R}(t)\|_{L^{3/2}}
   + |B_R|^{2/3}\|c_{k,R}-a_k\|_{L^{3/2}_t},
\]
which proves the asserted estimate for the fixed gauge \(a_k\).
\end{proof}

\begin{lemma}[Time derivative bound]\label{lem:time-derivative}
For every compact cylinder \(Q\Subset\R^3\times(-\infty,0)\), the sequence
\(\partial_tu_k\) is bounded in
\[
   L^{3/2}_tW^{-1,3/2}_x(Q).
\]
\end{lemma}

\begin{proof}
From the equation,
\[
   \partial_tu_k
   =
   \Delta u_k-\divergence(u_k\otimes u_k)-\nabla (p_k-a_k(t)),
\]
because \(\nabla a_k(t)=0\).  Let \(Q=B_R\times[-S,-\sigma]\).  The term
\(\Delta u_k\) is bounded in \(L^2_tH^{-1}_x(Q)\) by
\cref{lem:energy-diss}: for test fields \(\psi\in L^2_tH^1_{0,x}(Q)\),
\[
   \left|\iint \nabla u_k:\nabla\psi\right|
   \le \|\nabla u_k\|_{L^2(Q)}\|\nabla\psi\|_{L^2(Q)}.
\]
Because \(W^{1,3}_0(B_R)\subset H^1_0(B_R)\), an \(H^{-1}(B_R)\)-bound implies
a \(W^{-1,3/2}(B_R)\)-bound.  Since the time interval has finite length, the
\(L^2_tH^{-1}_x\)-bound is also an
\(L^{3/2}_tW^{-1,3/2}_x\)-bound.
The nonlinear term is bounded in \(L^{3/2}_tW^{-1,3/2}_x(Q)\) because, for
every \(\zeta\in W^{1,3}_0(B_R)\),
\[
   \left|\int_{B_R}(u_k\otimes u_k):\nabla\zeta\,dx\right|
   \le
   \|u_k\otimes u_k\|_{L^{3/2}(B_R)}
   \|\nabla\zeta\|_{L^3(B_R)} ,
\]
and
\[
   \|u_k\otimes u_k\|_{L^{3/2}(Q)}
   \le
   \|u_k\|_{L^\infty(Q)}
   \|u_k\|_{L^{3/2}(Q)}
   \le C(Q),
\]
using \cref{lem:inherited-bound}, the local \(L^2\)-bound, and the finite
measure of \(Q\).  Finally, for the pressure term, the dual definition of
\(W^{-1,3/2}\) gives, for every \(\zeta\in W^{1,3}_0(B_R)\),
\[
   \left|\int_{B_R}(p_k-a_k(t))\,\divergence \zeta\,dx\right|
   \le
   \|p_k-a_k(t)\|_{L^{3/2}(B_R)}
   \|\nabla\zeta\|_{L^3(B_R)} ,
\]
and hence
\[
   \|\nabla(p_k-a_k(t))\|_{W^{-1,3/2}(B_R)}
   \le \|p_k-a_k(t)\|_{L^{3/2}(B_R)},
\]
and \cref{lem:pressure} gives the required time integrability.
\end{proof}

\begin{proposition}[Compactness]\label{prop:compactness}
After passing to a subsequence,
\[
   u_k\to U
\]
strongly in \(L^q_{\loc}(\R^3\times(-\infty,0))\) for every finite \(q\), and
\[
   p_k-a_k(t)\rightharpoonup P
   \quad\hbox{weakly in }L^{3/2}_{\loc}(\R^3\times(-\infty,0)).
\]
\end{proposition}

\begin{proof}
Fix \(Q\Subset\R^3\times(-\infty,0)\).  By \cref{lem:energy-diss}, the
velocities are bounded in \(L^2_tH^1_x(Q)\).  By
\cref{lem:time-derivative}, the time derivatives are bounded in
\(L^{3/2}_tW^{-1,3/2}_x(Q)\).  On bounded spatial balls,
\[
   H^1\Subset L^2\hookrightarrow W^{-1,3/2}.
\]
The Aubin--Lions--Simon theorem \cite{Aubin1963,LionsMagenes,Simon1987}
therefore gives relative compactness in \(L^2(Q)\).  After passing to a
subsequence,
\(u_k\to U\) strongly in \(L^2(Q)\) and almost everywhere on \(Q\).

The inherited Type I bound gives a uniform \(L^\infty(Q)\) bound.  Interpolating
between strong \(L^2(Q)\) convergence and the common \(L^\infty(Q)\) bound
gives strong convergence in \(L^q(Q)\) for every finite \(q\).  The limit
belongs to the same \(L^\infty(Q)\) ball by weak-star lower semicontinuity in
\(L^\infty\), or equivalently by the pointwise a.e. convergence after passing
to a subsequence.  Thus, for \(q\ge2\),
\[
   \|u_k-U\|_{L^q(Q)}^q
   \le (2\sup_j\|u_j\|_{L^\infty(Q)}+2\|U\|_{L^\infty(Q)})^{q-2}
      \|u_k-U\|_{L^2(Q)}^2,
\]
and the case \(1\le q<2\) follows from Hölder on the finite-measure set
\(Q\).  Applying this argument to a countable exhaustion
\(Q_1\Subset Q_2\Subset\cdots\) of
\(\R^3\times(-\infty,0)\) by compact cylinders and selecting subsequences
successively gives a diagonal subsequence for which the convergence holds on
every \(Q_m\), and therefore locally.

Finally, \cref{lem:pressure} gives uniform local \(L^{3/2}\) bounds for the
gauged pressures \(p_k-a_k(t)\).  Since \(L^{3/2}\) is reflexive, the
Banach--Alaoglu theorem and the same diagonal extraction give a weak
\(L^{3/2}_{\mathrm{loc}}\) pressure limit.
\end{proof}

\section{The ancient suitable weak limit}\label{sec:compactness}

\begin{lemma}[Limit equation]\label{lem:limit-equation}
The limit \((U,P)\) obtained in \cref{prop:compactness} solves
\[
   \partial_tU+(U\cdot\nabla)U+\nabla P=\Delta U,\qquad
   \divergence U=0
\]
in distributions on \(\R^3\times(-\infty,0)\).
\end{lemma}

\begin{proof}
Let \(\varphi\in C_c^\infty(\R^3\times(-\infty,0);\R^3)\).  The weak
formulation for \((u_k,p_k-a_k(t))\) is
\[
   \iint \left[
      -u_k\cdot\partial_t\varphi
      -(u_k\otimes u_k):\nabla\varphi
      +(p_k-a_k(t))\divergence\varphi
      +\nabla u_k:\nabla\varphi
   \right]\,dx\,dt=0.
\]
The linear term with \(\partial_t\varphi\) passes to the limit by strong
\(L^2_{\loc}\)-convergence of \(u_k\).  The pressure term passes to the limit
by weak \(L^{3/2}_{\loc}\)-convergence of \(p_k-a_k(t)\), since
\(\divergence\varphi\in L^3\).  The viscous term passes to the limit by weak
\(L^2_{\loc}\)-convergence of \(\nabla u_k\).  The nonlinear term converges
because \(u_k\to U\) strongly in \(L^2_{\loc}\), hence
\[
   \|u_k\otimes u_k-U\otimes U\|_{L^1(\supp\varphi)}
   \le
   \|u_k-U\|_{L^2}
   \bigl(\|u_k\|_{L^2}+\|U\|_{L^2}\bigr)\to0 .
\]
Testing with gradients of scalar functions gives
\(\divergence U=0\), since \(\divergence u_k=0\) and
\(u_k\rightharpoonup U\) in \(L^2_{\loc}\).
\end{proof}

\begin{lemma}[Stability of suitability]\label{lem:suitability}
The pair \((U,P)\) satisfies the local energy inequality.  Hence it is a
suitable weak solution on \(\R^3\times(-\infty,0)\).
\end{lemma}

\begin{proof}
We use the distributional form of the local energy inequality.  Namely, for
every non-negative \(\phi\in C_c^\infty(\R^3\times(-\infty,0))\),
\[
   2\iint |\nabla u_k|^2\phi
   \le
   \iint |u_k|^2(\partial_t\phi+\Delta\phi)
   +\iint (|u_k|^2+2P_k)u_k\cdot\nabla\phi ,
\]
where \(P_k=p_k-a_k(t)\).  This formulation is obtained from
\eqref{eq:lei} by applying the endpoint inequality to time cutoffs and then
approximating the temporal cutoff monotonically; conversely it implies the
endpoint formulation for the weakly \(L^2_{\loc}\)-continuous representative.
Thus no additional time-slice convergence is needed in the compactness
argument.

The terms involving \(|u_k|^2(\partial_t\phi+\Delta\phi)\) converge strongly
in \(L^1\) because \(u_k\to U\) strongly in \(L^2_{\loc}\).  The cubic flux
term converges strongly in \(L^1\) because \(u_k\to U\) strongly in
\(L^3_{\loc}\).

For the pressure flux, since \(P_k\rightharpoonup P\) in \(L^{3/2}\) and
\[
   u_k\cdot\nabla\phi\to U\cdot\nabla\phi
   \quad\hbox{strongly in }L^3,
\]
the duality pairing
\(\iint P_k u_k\cdot\nabla\phi\) converges to
\(\iint P U\cdot\nabla\phi\).  Finally,
\(\nabla u_k\rightharpoonup\nabla U\) in \(L^2_{\loc}\), so the convex
functional
\[
   w\mapsto \iint |w|^2\phi
\]
is weakly lower semicontinuous.  Taking the lower limit on the left-hand side
and the ordinary limits on the right-hand side gives the local energy
inequality for \(U\) in distributional form.  The endpoint form in
\cref{def:suitable} follows by the same monotone temporal-cutoff
approximation and the standard weak \(L^2_{\loc}\)-continuity of local
energy-class weak solutions.  The distributional equation, divergence-free
condition, and local integrability requirements were already obtained in
\cref{lem:limit-equation,prop:compactness}; hence \((U,P)\) is suitable.
\end{proof}

\begin{lemma}[Inherited ancient Type I bound]\label{lem:limit-bound}
For every compact interval \(I=[-S,-\sigma]\Subset(-\infty,0)\),
\[
   \|U\|_{L^\infty(\R^3\times I)}
   \le M\sigma^{-1/2}.
\]
\end{lemma}

\begin{proof}
Let \(I=[-S,-\sigma]\Subset(-\infty,0)\).  By
\cref{lem:inherited-bound},
\[
   \|u_k\|_{L^\infty(\R^3\times I)}
   \le M\sigma^{-1/2}.
\]
Fix \(R<\infty\).  The convergence in \cref{prop:compactness} gives, after
passing to a subsequence if necessary, \(u_k\to U\) a.e. on \(B_R\times I\).
If a measurable set \(E\subset B_R\times I\) had positive measure and
\(|U|>M\sigma^{-1/2}+\delta\) on \(E\), then by a.e. convergence the functions
\(|u_k|\) would exceed \(M\sigma^{-1/2}+\delta/2\) on a subset of \(E\) of
positive measure for all sufficiently large \(k\), contradicting the uniform
bound for \(u_k\).  Hence
\[
   \|U\|_{L^\infty(B_R\times I)}\le M\sigma^{-1/2}.
\]
Letting \(R\to\infty\) through integers gives the asserted bound on
\(\R^3\times I\).
\end{proof}

\begin{proposition}[Ancient suitable Type I limit]\label{prop:ancient-limit}
The pair \((U,P)\) is an ancient suitable weak solution on
\(\R^3\times(-\infty,0)\) satisfying the interval Type I bound in
\cref{lem:limit-bound}.
\end{proposition}

\begin{proof}
Combine \cref{lem:limit-equation,lem:suitability,lem:limit-bound}.
\end{proof}

\section{Non-triviality}\label{sec:nontrivial}

We now choose the scales so that the ancient limit is not zero.  This is the
only selection step in the proof, and we separate the local regularity input
from the scale selection.

\begin{lemma}[Pressure-normalized concentration at a singular point]
\label{lem:singular-concentration}
If \(z_*=(x_*,T)\) is singular, then there are \(\eta>0\) and \(r_0>0\) such
that, for every \(0<r<r_0\),
\begin{equation}\label{eq:positive-concentration-sequence}
   C(u;z_*,r)+D(p;z_*,r)\ge \eta .
\end{equation}
\end{lemma}

\begin{proof}
Take \(\eta=\eps_*\), where \(\eps_*\) is the epsilon-regularity constant in
\cref{thm:eps-reg}.  If there were arbitrarily small radii \(r\) with
\[
   C(u;z_*,r)+D(p;z_*,r)<\eps_*,
\]
then \cref{thm:eps-reg} would imply that \(u\) is bounded in
\(Q_{r/2}(z_*)\).  This is precisely regularity of \(z_*\) in
\cref{def:terminal-singular}, contradicting the hypothesis that \(z_*\) is
singular.  Hence the lower bound holds for all sufficiently small \(r\).
\end{proof}

\begin{lemma}[Endpoint compactness at a centered Type I point]
\label{lem:endpoint-compactness}
Assume \(\calI_*(u,p;z_*)<\infty\).  Let \(\lambda_k\downarrow0\), and let
\((u_k,p_k)\) be the centered rescalings \eqref{eq:rescaled-sequence}.  For
every \(R<\infty\) there are functions of time \(b_{k,R}(t)\) such that,
after passing to a subsequence,
\[
   u_k\to U
   \quad\hbox{strongly in }L^3(Q_{R'})
   \quad\hbox{for every }0<R'<R,
 \]
and
\[
   p_k-b_{k,R}(t)\rightharpoonup P
   \quad\hbox{weakly in }L^{3/2}(Q_R).
 \]
The limit is suitable in \(Q_R\).  A diagonal subsequence gives this conclusion
for every finite \(R\).
\end{lemma}

\begin{proof}
By the scale invariance in \cref{lem:scaling}, the hypothesis
\(\calI_*(u,p;z_*)<\infty\) gives, for each fixed \(R\),
\[
   \sup_k\left[
   A(u_k;0,R)+C(u_k;0,R)+D(p_k;0,R)+E(u_k;0,R)
   \right]<\infty
\]
once \(k\) is sufficiently large.  The pressure term is understood after
subtracting the spatial average on \(B_R\), which is the function \(b_{k,R}\).
The local energy inequality is invariant under this pressure gauge, since the
subtracted term is a function of time and \(u_k\) is divergence free.

We spell out the compactness.  The \(A\)-bound gives
\[
   \sup_k\|u_k\|_{L^\infty(-R^2,0;L^2(B_R))}<\infty,
\]
and the \(E\)-bound gives
\[
   \sup_k\|\nabla u_k\|_{L^2(Q_R)}<\infty.
\]
By the Sobolev embedding \(H^1(B_R)\hookrightarrow L^6(B_R)\), the
\(E\)-bound gives \(u_k\in L^2_tL^6_x(Q_R)\).  Interpolating
\[
   \frac{1}{10/3}=\frac25\cdot\frac12+\frac35\cdot\frac16
\]
on each time slice gives
\[
   \|u_k(t)\|_{L^{10/3}(B_R)}
   \le
   \|u_k(t)\|_{L^2(B_R)}^{2/5}
   \|u_k(t)\|_{L^6(B_R)}^{3/5}.
\]
Raising to the power \(10/3\) and integrating in time yields
\[
   \int_{-R^2}^0\|u_k(t)\|_{L^{10/3}(B_R)}^{10/3}\,dt
   \le
   \left(\sup_{-R^2<t<0}\|u_k(t)\|_{L^2(B_R)}^2\right)^{2/3}
   \int_{-R^2}^0\|u_k(t)\|_{L^6(B_R)}^2\,dt .
\]
Thus the \(A\)- and \(E\)-bounds imply the energy-class estimate
\[
   \sup_k\|u_k\|_{L^{10/3}(Q_R)}<\infty.
\]
The equation written with the pressure gauge \(p_k-b_{k,R}(t)\) gives
\[
   \partial_tu_k
   =
   \Delta u_k-\divergence(u_k\otimes u_k)
   -\nabla(p_k-b_{k,R}(t)).
\]
The three terms on the right are bounded in negative Sobolev spaces on every
strictly smaller cylinder.  The term \(\Delta u_k\) is bounded in
\(L^2_tH^{-1}_x\), hence in
\(L^{3/2}_tW^{-1,3/2}_x\) on finite time intervals, by the \(E\)-bound and
the embedding \(W^{1,3}_0(B_{R'})\subset H^1_0(B_{R'})\).  Since
\(u_k\in L^{10/3}(Q_R)\),
\[
   u_k\otimes u_k\in L^{5/3}(Q_R)\subset L^{3/2}(Q_{R'})
\]
on the finite-measure cylinder \(Q_{R'}\), and therefore
\(\divergence(u_k\otimes u_k)\) is bounded in
\(L^{3/2}_tW^{-1,3/2}_x(Q_{R'})\).  The \(D\)-bound gives
\[
   p_k-b_{k,R}(t)\in L^{3/2}(Q_R),
\]
so \(\nabla(p_k-b_{k,R}(t))\) is bounded in the same negative Sobolev space.
Hence, for each \(R'<R\),
\(\partial_tu_k\) is bounded in
\[
   L^{3/2}((-R'^2,0);W^{-1,3/2}(B_{R'})).
\]
Since \(H^1(B_{R'})\Subset L^2(B_{R'})\hookrightarrow W^{-1,3/2}(B_{R'})\),
the Aubin--Lions--Simon theorem \cite{Aubin1963,LionsMagenes,Simon1987} gives
strong convergence in \(L^2(Q_{R'})\).
The \(L^{10/3}\)-bound then upgrades this to strong convergence in
\(L^3(Q_{R'})\) by interpolation:
\[
   \|u_k-U\|_{L^3(Q_{R'})}
   \le
   \|u_k-U\|_{L^2(Q_{R'})}^{\theta}
   \|u_k-U\|_{L^{10/3}(Q_{R'})}^{1-\theta}
   \to0
\]
for the unique \(\theta\in(0,1)\) satisfying
\(1/3=\theta/2+(1-\theta)/(10/3)\).  The \(D\)-bound is precisely the
\(L^{3/2}(Q_R)\)-bound for \(p_k-b_{k,R}(t)\), so reflexivity gives a weak
\(L^{3/2}\)-limit.

It remains to record why the limiting pair is suitable in \(Q_R\).  Let
\(\phi\ge0\) be compactly supported in \(Q_R\).  Apply the local energy
inequality to \((u_k,p_k-b_{k,R}(t))\).  The terms with
\(|u_k|^2(\partial_t\phi+\Delta\phi)\) converge because
\(u_k\to U\) strongly in \(L^2\); the cubic flux converges because
\(u_k\to U\) strongly in \(L^3\); and the pressure flux converges by the
duality between weak \(L^{3/2}\)-convergence of \(p_k-b_{k,R}(t)\) and strong
\(L^3\)-convergence of \(u_k\cdot\nabla\phi\).  The dissipation term is lower
semicontinuous under weak \(L^2\)-convergence of \(\nabla u_k\).  Thus the
local energy inequality passes to the limit.  The distributional equation and
the divergence-free constraint pass to the limit by the same strong--weak
convergences used above.  Hence the limit is suitable in \(Q_R\).  A diagonal
subsequence over \(R=1,2,\ldots\) gives the assertion for all finite \(R\).
\end{proof}

\begin{lemma}[Scale selection from local concentration]\label{lem:scale-selection}
The scales in \cref{thm:main} may be chosen from a positive local
concentration sequence.  More precisely, there are
\(\lambda_k\downarrow0\), a compact cylinder
\[
   K_0\Subset \R^3\times(-\infty,0),
\]
the reference gauges \(a_k(t)=(p_k)_{B_1}(t)\) from
\cref{lem:pressure}, and a number \(c_0>0\) such that the centered
rescalings \eqref{eq:rescaled-sequence} satisfy
\begin{equation}\label{eq:selected-lower-bound}
   \iint_{K_0}
   \left(|u_k|^3+|p_k-a_k(t)|^{3/2}\right)\,dx\,dt\ge c_0
\end{equation}
for all \(k\).  In addition, the compactness limit obtained from this selected
sequence is nonzero in \(L^3_{\loc}(\R^3\times(-\infty,0))\).
\end{lemma}

\begin{proof}
By \cref{lem:singular-concentration}, every sufficiently small radius carries
positive pressure-normalized density.  Choose any sequence \(r_j\downarrow0\)
inside that range, so that
\[
   C(u;z_*,r_j)+D(p;z_*,r_j)\ge \eta .
\]
After rescaling by \(r_j\), this says that the pressure-normalized density of
the pair
\[
   v_j(x,t)=r_j u(x_*+r_jx,T+r_j^2t),\qquad
   q_j(x,t)=r_j^2p(x_*+r_jx,T+r_j^2t)
\]
on \(Q_1\) is bounded below by \(\eta\), with \(q_j\) normalized by subtracting
\((q_j)_{B_1}(t)\).  The subtraction is a legitimate pressure gauge by
\cref{lem:gauges-harmless}.

We now invoke the Caffarelli--Kohn--Nirenberg--Seregin concentration
compactness selection for singular points
\cite{CKN1982,Lin1998,Seregin2012}.  In the local Type I setting this is the
forward implication in Albritton--Barker
\cite[Theorem 1.1]{AlbrittonBarker2019}: a positive
pressure-normalized concentration sequence at a Type I singular terminal point
has a subsequence, possibly after replacing the radii by comparable radii
\(\lambda_j\downarrow0\), whose rescalings converge to a nonzero bounded
ancient limit.  The pressures in that compactness statement are normalized by
a fixed spatial average, and after the whole-space representative has been
chosen as in \cref{lem:pressure} we take this average to be the reference
gauge \(a_j(t)=(p_j)_{B_1}(t)\).  This is a gauge choice only, by
\cref{lem:gauges-harmless}, and it is compatible with the local pressure
compactness used in \cref{prop:compactness}.  In the notation used here, the
nonzero conclusion is precisely that a fixed compact subcylinder
\(K_0\Subset Q_1\cap\{t<0\}\) carries a uniform positive amount of
pressure-normalized density.  Thus there are \(K_0\) and \(c_0>0\) such that
\[
   \iint_{K_0}
   \left(|u_j|^3+|p_j-a_j(t)|^{3/2}\right)\,dx\,dt\ge c_0 .
\]
Here \(u_j,p_j\) denote the centered rescalings at the selected scales
\(\lambda_j\), and \(a_j(t)=(p_j)_{B_1}(t)\).

No additional hypothesis is being imposed in this selection step.  The input is
only the positive CKN concentration supplied by
\cref{lem:singular-concentration}, the local Type I alternative in the
definition of the source point, and the endpoint compactness supplied by
\(\calI_*(u,p;z_*)<\infty\).  The possible loss of density into the terminal
slice is the defect excluded by the cited Seregin selection theorem.

The global Type I bound \eqref{eq:type-I-bound-intro} rules out the Type II
branch of the local blow-up-rate dichotomy along the selected sequence, while
\(\calI_*(u,p;z_*)<\infty\) gives the endpoint compactness needed above.
The selected sequence is therefore a sequence for which the compactness limit
is nonzero.  Renaming that selected sequence as \(k\) gives
\eqref{eq:selected-lower-bound} and the final assertion of the lemma.
\end{proof}

\begin{remark}\label{rem:selection}
The scale-selection lemma is the only point at which singularity, rather than
mere bounded Type I behavior, is used.  It is also the point at which the
terminal-slice concentration defect is excluded.  For this reason the lower
bound \eqref{eq:selected-lower-bound} is placed on a cylinder compactly
contained in \(\{t<0\}\).
\end{remark}

\begin{proposition}[Nonzero ancient limit]\label{prop:nonzero}
For the scales chosen in \cref{lem:scale-selection}, the limit \(U\) in
\cref{prop:ancient-limit} is nonzero.  More precisely, there are a compact
cylinder \(K_0\Subset\R^3\times(-\infty,0)\), a time-dependent pressure gauge
\(a(t)\), and a number \(\eta_1>0\) such that
\[
   \iint_{K_0}\left(|U|^3+|P-a(t)|^{3/2}\right)\,dx\,dt\ge \eta_1 .
\]
\end{proposition}

\begin{proof}
The scale-selection theorem invoked in \cref{lem:scale-selection} is used here
with its nonzero-limit conclusion.  We apply the compactness construction of
\cref{prop:compactness} to that already selected sequence.  If a further
subsequence is needed for local compactness, it is taken from a sequence which
already converges, in the cited selection theorem, to a bounded ancient
suitable weak solution which is not identically zero.  The further subsequence
therefore has the same limit.  Thus the limit denoted by \(U\) in
\cref{prop:ancient-limit} satisfies
\[
   U\not\equiv0
   \quad\hbox{in }L^3_{\loc}(\R^3\times(-\infty,0)).
\]

Since \(U\in L^3_{\loc}\) and is not zero, there is a compact cylinder
\(K_0\Subset\R^3\times(-\infty,0)\) such that
\[
   \iint_{K_0}|U|^3\,dx\,dt>0 .
\]
Let \(a(t)\) be any time-dependent pressure gauge for \(P\) on \(K_0\), for
instance the spatial average of \(P\) over a ball containing the spatial
projection of \(K_0\).  Then
\[
   \iint_{K_0}\left(|U|^3+|P-a(t)|^{3/2}\right)\,dx\,dt
   \ge
   \iint_{K_0}|U|^3\,dx\,dt .
\]
Taking \(\eta_1\) to be the positive right-hand side gives the asserted
pressure-normalized lower bound.  No lower semicontinuity of the pressure norm
is used in this step.
\end{proof}

\section{Centered self-similar variables}\label{sec:renorm}

The Type I bound makes the following change of variables natural:
\[
   y=\frac{x}{\sqrt{-t}},\qquad
   \tau=-\log(-t),\qquad
   V(y,\tau)=\sqrt{-t}\,U(y\sqrt{-t},t).
\]
Set also
\[
   \Pi(y,\tau)=(-t)P(y\sqrt{-t},t),
\]
up to addition of a function of \(\tau\).

\begin{lemma}[Renormalized equation]\label{lem:renormalized-equation}
The pair \((V,\Pi)\) solves
\[
   \partial_\tau V+(V\cdot\nabla)V+\nabla\Pi
   =
   \Delta V-\frac12(V+y\cdot\nabla V),
   \qquad
   \divergence V=0.
\]
\end{lemma}

\begin{proof}
Write \(s=-t=e^{-\tau}\), \(x=s^{1/2}y\), and
\[
   U(x,t)=s^{-1/2}V(y,\tau).
\]
The identities below are identities in distributions obtained by testing
against the inverse change of variables.  Since \cref{lem:smoothness} gives
smooth representatives on compact subsets of \(t<0\), the same computation is
classical for the representative appearing in the theorem.
Since \(s_t=-1\), \(\tau_t=s^{-1}\), and \(y_t=(2s)^{-1}y\), one computes
\[
   \partial_tU
   =
   s^{-3/2}\left(\partial_\tau V+\frac12V+\frac12y\cdot\nabla V\right),
\]
\[
   (U\cdot\nabla_x)U=s^{-3/2}(V\cdot\nabla)V,\qquad
   \Delta_xU=s^{-3/2}\Delta V,\qquad
   \nabla_xP=s^{-3/2}\nabla\Pi.
\]
Substitution in the Navier--Stokes equation for \(U\) gives the displayed
equation.  The divergence condition is preserved by the spatial scaling.
\end{proof}

\begin{lemma}[Renormalized boundedness and non-triviality]
\label{lem:renorm-bounds}
The field \(V\) satisfies
\[
   \|V\|_{L^\infty(\R^3\times\R)}\le M.
\]
Moreover \eqref{eq:main-nonzero-renorm} holds for some compact cylinder
\(K\Subset\R^3\times\R\), function \(b(\tau)\), and \(\eta_0>0\).
\end{lemma}

\begin{proof}
The \(L^\infty\) bound follows from the classical pointwise Type I bound in
\cref{cor:classical-type-I-bound}:
\[
   |V(y,\tau)|
   =
   \sqrt{-t}\,|U(y\sqrt{-t},t)|
   \le M.
\]
For non-triviality, \cref{prop:nonzero} gives positive pressure-normalized CKN
density on a compact spacetime cylinder in \((x,t)\).  Under the change of
variables \((x,t)\mapsto(y,\tau)\), this compact cylinder maps into a compact
set contained in some compact cylinder \(K\Subset\R^3\times\R\).  If
\(s=-t=e^{-\tau}\), then
\[
   dx\,dt=s^{5/2}\,dy\,d\tau,\qquad
   U=s^{-1/2}V,\qquad
   P-a(t)=s^{-1}\bigl(\Pi-b(\tau)\bigr),
\]
where \(b(\tau)=s\,a(-s)\).  Therefore
\[
   |U|^3\,dx\,dt
   =
   s\,|V|^3\,dy\,d\tau,
   \qquad
   |P-a(t)|^{3/2}\,dx\,dt
   =
   s\,|\Pi-b(\tau)|^{3/2}\,dy\,d\tau .
\]
On the image of the compact cylinder, \(s\) is bounded above and below by
positive constants.  Thus a positive lower bound in \((x,t)\)-variables is
equivalent, up to a fixed multiplicative constant, to a positive lower bound
in \((y,\tau)\)-variables.  This gives \eqref{eq:main-nonzero-renorm} with a
possibly different positive constant \(\eta_0\).
\end{proof}

\begin{lemma}[Smoothness]\label{lem:smoothness}
The ancient solution \(U\) is smooth on
\(\R^3\times(-\infty,0)\), and \(V\) is smooth on \(\R^3\times\R\).
\end{lemma}

\begin{proof}
On each compact cylinder \(Q\Subset\R^3\times(-\infty,0)\), the bound
\cref{lem:limit-bound} gives \(U\in L^\infty(Q)\).  In particular
\(U\in L^s_tL^r_x(Q)\) for every finite \(r,s\).  Choose, for example,
\(r=s=10\); then \(2/s+3/r=1/2<1\).  Serrin's interior regularity criterion
\cite{Serrin1962,Sohr2001} for suitable weak solutions applies on \(Q\).  In
the form used here, the criterion asserts that every point of \(Q\) is regular:
for each \(Q'\Subset Q\), \(U\) is bounded and H\"older continuous on \(Q'\),
and the interior bootstrap for the Stokes system gives smoothness on still
smaller cylinders.  The bootstrap starts from the equation
\[
   \partial_tU-\Delta U+\nabla P=-\divergence(U\otimes U),
   \qquad \divergence U=0,
\]
with \(U\otimes U\) locally bounded after Serrin regularity.  Interior Stokes
estimates on nested cylinders give higher spatial derivatives and one time
derivative; differentiating the equation and iterating the same estimates gives
\[
   U\in C^\infty(\R^3\times(-\infty,0)).
\]
The pressure is recovered locally from
\(-\Delta P=\partial_i\partial_j(U_iU_j)\) plus a harmonic function and is
smooth after fixing a time-dependent gauge.  The self-similar change of
variables is a smooth diffeomorphism for \(t<0\), so \(V\) and \(\Pi\) are
smooth on \(\R^3\times\R\).
\end{proof}

\begin{corollary}[Classical Type I bound]\label{cor:classical-type-I-bound}
For the smooth representative from \cref{lem:smoothness},
\[
   \|U(t)\|_{L^\infty(\R^3)}
   \le \frac{M}{\sqrt{-t}},
   \qquad t<0.
\]
\end{corollary}

\begin{proof}
Fix \(t_0<0\).  Choose \(\delta_n\downarrow0\) such that
\[
   I_n=[t_0-\delta_n,t_0+\delta_n]\Subset(-\infty,0).
\]
By \cref{lem:limit-bound},
\[
   \|U\|_{L^\infty(\R^3\times I_n)}
   \le M(-t_0-\delta_n)^{-1/2}.
\]
Since the representative is continuous in time and space, this essential
spacetime bound holds pointwise on \(\R^3\times I_n\): otherwise a point where
the bound is exceeded would have a small spacetime neighborhood of positive
measure where it is still exceeded.  Hence, for every \(x\in\R^3\),
\[
   |U(x,t_0)|
   \le M(-t_0-\delta_n)^{-1/2}
   \qquad\hbox{for every }n.
\]
Letting \(n\to\infty\) and then taking the supremum over \(x\in\R^3\) gives
the asserted estimate.
\end{proof}

\begin{proof}[Proof of \cref{thm:main}]
The compact ancient limit and the interval Type I bound are
\cref{prop:ancient-limit}; the pointwise bound
\eqref{eq:ancient-type-I-bound} is \cref{cor:classical-type-I-bound}.
Non-triviality is \cref{prop:nonzero}.  The renormalized equation is
\cref{lem:renormalized-equation}, and the boundedness and renormalized
non-triviality are \cref{lem:renorm-bounds}.  Smoothness is
\cref{lem:smoothness}.
\end{proof}

\section{Absence of additional centered symmetries}\label{sec:symmetry}

\begin{lemma}[Residual logarithmic-time translation]\label{lem:tau-translation}
If \(V\) solves \eqref{eq:centered-intro}, then for every \(\tau_0\in\R\)
the function
\[
   V_{\tau_0}(y,\tau)=V(y,\tau+\tau_0)
\]
also solves \eqref{eq:centered-intro}.  Moreover, if
\[
   U^\lambda(x,t)=\lambda U(\lambda x,\lambda^2t)
\]
is a parabolic rescaling of the unrenormalized ancient solution, then the
renormalized field associated with \(U^\lambda\) is
\[
   V^\lambda(y,\tau)=V(y,\tau-2\log\lambda).
\]
\end{lemma}

\begin{proof}
The first statement follows because \eqref{eq:centered-intro} is autonomous in
\(\tau\).  For the second, write \(s=-t=e^{-\tau}\).  Then
\[
   \sqrt{s}\,U^\lambda(y\sqrt{s},-s)
   =
   \lambda\sqrt{s}\,U(\lambda y\sqrt{s},-\lambda^2s)
   =
   V(y,-\log(\lambda^2s))
   =
   V(y,\tau-2\log\lambda).
\]
\end{proof}

\begin{lemma}[Failure of scaling invariance]\label{lem:no-scaling}
Let \(V\) be a bounded smooth solution of \eqref{eq:centered-intro}.  For
\(\alpha>0\), define
\[
   \widetilde V(y,\tau)=\alpha V(\alpha y,\alpha^2\tau).
\]
If \(V\not\equiv0\) and \(\widetilde V\) solves the same centered equation
with the pressure
\(\widetilde\Pi(y,\tau)=\alpha^2\Pi(\alpha y,\alpha^2\tau)\), then
\(\alpha=1\).
\end{lemma}

\begin{proof}
Write \(Y=\alpha y\) and \(T=\alpha^2\tau\).  The terms
\[
   \partial_\tau V,\qquad (V\cdot\nabla)V,\qquad \nabla\Pi,\qquad \Delta V
\]
in the non-drift part scale with the common factor \(\alpha^3\).  More
precisely,
\[
   \partial_\tau\widetilde V
   +(\widetilde V\cdot\nabla)\widetilde V
   +\nabla\widetilde\Pi-\Delta\widetilde V
   =
   \alpha^3
   \left(\partial_TV+(V\cdot\nabla_Y)V+\nabla_Y\Pi-\Delta_YV\right)(Y,T).
\]
The drift term satisfies
\[
   -\frac12(\widetilde V+y\cdot\nabla\widetilde V)
   =
   -\frac{\alpha}{2}
   \bigl(V+Y\cdot\nabla_Y V\bigr)(Y,T).
\]
Using the equation for \(V\), the residual of the centered equation for
\(\widetilde V\) is therefore
\[
   \frac{\alpha-\alpha^3}{2}
   \bigl(V+Y\cdot\nabla_YV\bigr)(Y,T).
\]
If \(\widetilde V\) solves the equation and \(\alpha\ne1\), then
\[
   V(Y,T)+Y\cdot\nabla_YV(Y,T)=0
   \qquad\hbox{for all }(Y,T).
\]
Fix \(T\) and a direction \(\omega\in S^2\).  For \(r>0\),
\[
   \frac{d}{dr}\bigl(rV(r\omega,T)\bigr)
   =
   V(r\omega,T)+r\,\omega\cdot\nabla V(r\omega,T)=0.
\]
Hence \(rV(r\omega,T)\) is constant in \(r\).  Since \(V\) is bounded and
smooth at \(Y=0\), this constant must be zero, and so \(V(r\omega,T)=0\) for
all \(r>0\).  By continuity, \(V(\cdot,T)\equiv0\), and since \(T\) was
arbitrary, \(V\equiv0\).  This contradicts \(V\not\equiv0\); hence
\(\alpha=1\).
\end{proof}

\begin{lemma}[Failure of translation invariance]\label{lem:no-translation}
Let \(a\in\R^3\), and define \(\widetilde V(y,\tau)=V(y-a,\tau)\).  The map
\(V\mapsto\widetilde V\) is a symmetry of the centered equation only when
\(a=0\).  For a fixed solution \(V\), the translated field can solve the same
equation only if \(a\cdot\nabla V\equiv0\).
\end{lemma}

\begin{proof}
All translation-invariant terms in the equation transform by composition with
the shifted variable.
The drift does not:
\[
   y\cdot\nabla\widetilde V(y,\tau)
   =
   y\cdot\nabla V(y-a,\tau)
   =
   (y-a)\cdot\nabla V(y-a,\tau)
   +a\cdot\nabla V(y-a,\tau).
\]
The extra term \(a\cdot\nabla V(y-a,\tau)\) is absent from the original
equation.  Hence translation invariance fails unless \(a=0\).  For a fixed
solution, the translated field can still solve the same equation only in the
exceptional case \(a\cdot\nabla V\equiv0\).
\end{proof}

\begin{proof}[Proof of \cref{prop:no-cascade}]
\Cref{lem:tau-translation} identifies the residual action of the original
parabolic scaling: it is translation in logarithmic time.  \Cref{lem:no-scaling}
excludes any additional non-trivial amplitude-space scaling symmetry, and
\cref{lem:no-translation} shows that constant spatial recentering is not a
symmetry of the fixed centered equation.  Therefore any multi-scale behavior of
a bounded ancient profile must be encoded in the evolution of that profile in
\(\tau\), not in an orbit generated by further centered rescalings.
\end{proof}

\begin{thebibliography}{99}

\bibitem{AlbrittonBarker2019}
D. Albritton and T. Barker,
\emph{On local Type I singularities of the Navier--Stokes equations and
Liouville theorems},
J. Math. Fluid Mech. \textbf{21} (2019), article no. 43,
doi:\url{https://doi.org/10.1007/s00021-019-0448-z}.

\bibitem{Aubin1963}
J.-P. Aubin,
\emph{Un theoreme de compacite},
C. R. Acad. Sci. Paris \textbf{256} (1963), 5042--5044.

\bibitem{CKN1982}
L. Caffarelli, R. Kohn, and L. Nirenberg,
\emph{Partial regularity of suitable weak solutions of the Navier--Stokes equations},
Comm. Pure Appl. Math. \textbf{35} (1982), no. 6, 771--831,
doi:\url{https://doi.org/10.1002/cpa.3160350604}.

\bibitem{ESS2003}
L. Escauriaza, G. Seregin, and V. \v{S}verak,
\emph{\(L_{3,\infty}\)-solutions of Navier--Stokes equations and backward uniqueness},
Russian Math. Surveys \textbf{58} (2003), no. 2, 211--250,
doi:\url{https://doi.org/10.1070/RM2003v058n02ABEH000609}.

\bibitem{Lin1998}
F.-H. Lin,
\emph{A new proof of the Caffarelli--Kohn--Nirenberg theorem},
Comm. Pure Appl. Math. \textbf{51} (1998), no. 3, 241--257,
doi:\url{https://doi.org/10.1002/(SICI)1097-0312(199803)51:3<241::AID-CPA2>3.0.CO;2-A}.

\bibitem{Leray1934}
J. Leray,
\emph{Sur le mouvement d'un liquide visqueux emplissant l'espace},
Acta Math. \textbf{63} (1934), 193--248.

\bibitem{LionsMagenes}
J.-L. Lions and E. Magenes,
\emph{Non-homogeneous boundary value problems and applications. Vol. I},
Springer-Verlag, 1972.

\bibitem{Seregin2012}
G. Seregin,
\emph{A certain necessary condition of potential blow up for Navier--Stokes equations},
Comm. Math. Phys. \textbf{312} (2012), no. 3, 833--845,
doi:\url{https://doi.org/10.1007/s00220-011-1391-x}.

\bibitem{Serrin1962}
J. Serrin,
\emph{On the interior regularity of weak solutions of the Navier--Stokes equations},
Arch. Rational Mech. Anal. \textbf{9} (1962), 187--195,
doi:\url{https://doi.org/10.1007/BF00253344}.

\bibitem{Simon1987}
J. Simon,
\emph{Compact sets in the space \(L^p(0,T;B)\)},
Ann. Mat. Pura Appl. (4) \textbf{146} (1987), 65--96,
doi:\url{https://doi.org/10.1007/BF01762360}.

\bibitem{Sohr2001}
H. Sohr,
\emph{The Navier--Stokes equations. An elementary functional analytic approach},
Birkhauser, 2001.

\bibitem{Stein1970}
E. M. Stein,
\emph{Singular integrals and differentiability properties of functions},
Princeton University Press, 1970.

\end{thebibliography}

\end{document}
